\newgeometry{top = 0.7cm, left = 1cm, right = 1cm, bottom=0.1cm} %zmenšení okrajů, aby se vešla na stránku
\song{Těšínská (Jaromír Nohavica)}

\thispagestyle{empty}

\vspace{-5px}
\ve{}
\ch{Ami}Kdybych se narodil \ch{Dmi}před sto léty \ch{F} \\
\ch{E7}v tomhle \ch{Ami}městě, \ch{Dmi F E7 Ami}\\
\ch{Ami}u Larischů na zahradě \ch{Dmi}trhal bych květy \ch{F}\\
\ch{E7}své ne\ch{Ami}věstě.\ch{Dmi F E7 Ami}\\[2px]
\ch{C}Moje nevěsta by \ch{Dmi}byla dcera ševcova\\
\ch{F}z domu Kamińskich \ch{C}odněkud ze Lvova\\
kochal bym ja i \ch{Dmi}pieśćił \ch{F}chy\ch{E7}ba lat \ch{Ami}dwieśćie. 

\vspace{-5px}
\ve{}
Bydleli bychom na Sachsenbergu\\
v domě u žida Kohna.\\
Nejhezčí ze všech těšínských šperků\\
byla by ona.\\[2px]
Mluvila by polsky a trochu česky,\\
pár slov německy a smála by se hezky.\\
Jednou za sto let zázrak se koná, zázrak se koná.

\vspace{-5px}
\ve{}
Kdybych se narodil před sto léty\\
byl bych vazačem knih.\\
U Prohazků dělal bych od pěti do pěti\\
a sedm zlatek za to bral bych.\\[2px]
Měl bych krásnou ženu a tři děti,\\
zdraví bych měl a bylo by mi kolem třiceti,\\
celý dlouhý život před sebou celé krásné dvacáté století.

\vspace{-5px}
\ve{}
Kdybych se narodil před sto léty,\\
v jinačí době,\\
u Larichů na zahradě trhal bych květy\\
má lásko tobě.\\[2px]
Tramvaj by jezdila přes řeku nahoru,\\
slunce by zvedalo hraniční závoru\\
a z oken voněl by sváteční oběd.

\vspace{-5px}
\ve{}
Večer by zněla od Mojzese\\
melodie dávnověká,\\
bylo by léto tisíc devět set deset\\
a za domem by tekla řeka.\\[2px]
Vidím to jako dnes šťastného sebe,\\
ženu a děti a těšínské nebe.\\
Jěštěže člověk nikdy neví co ho čeká.\\[2px]
na na na na... 

\esong{}
\restoregeometry{}